\documentclass[12pt]{article}
\usepackage[utf8]{inputenc}
\usepackage[T1]{fontenc}
\usepackage{hyperref}
\title{Kuuk Thayorre Language Resources for K-12 Curriculum Development}
\author{Compiled by Claude Sonnet 5 (claude-sonnet-5) for Andr\'as Kornai}
\date{}
\begin{document}
\maketitle
\section{Introduction}
Kuuk Thaayorre (ISO 639-3: thd) is a Paman (Pama-Nyungan) language of Cape York, Queensland, associated with Pormpuraaw. Wikipedia lists 174 speakers (2021 census) and UNESCO's ``Definitely Endangered'' status. The searches did not turn up a dictionary or a formal school curriculum.
\section{Primary Educational Resources}
A BYU Stowaway blog post says the language is still being taught to Thaayorre children, without describing a formal curriculum: \url{https://stowawaymag.byu.edu/blog/kuuk-thaayorre-an-endangered-language-moving-in-the-right-direction}
\section{Linguistic and Reference Materials}
\begin{itemize}
\item Alice R. Gaby, A Grammar of Kuuk Thaayorre (De Gruyter Mouton, 2017; Mouton Grammar Library 74; paperback 2019, ISBN 9783110653304; fieldwork 2002--2008): \url{https://www.sack.de/gaby-a-grammar-of-kuuk-thaayorre/9783110653304}
\item LINGUIST List review of the Gaby grammar: \url{https://linguistlist.org/issues/30/33/}
\item Wikipedia, Kuuk Thaayorre language: \url{https://en.wikipedia.org/wiki/Kuuk_Thaayorre_language}
\item Omniglot page with alphabet chart: \url{https://omniglot.com/writing/kuukthaayorre.htm}
\end{itemize}
\section{Religious and Community Resources}
None found.
\section{Digital and Audio Resources}
None found.
\section{Curriculum Development Context}
A technical grammar exists but is not a classroom resource. Next steps would be to contact the Pormpuraaw community and local language programme directly (not verified here), and check AIATSIS for dictionaries and recordings.
\end{document}